% Section 5.4 — The Unassociated Source Problem
% Trimmed: condensed association completeness, compressed previous search chronology.
% Absorbed 5.5 concepts: quantification learning motivation as closing paragraphs.

\section{The Unassociated Source Problem}
\label{sec:5.4}

Having established the theoretical grounds for dark matter subhalo searches in the previous sections, we now turn to the observational reality.
The Fermi-LAT catalogs contain a large population of gamma-ray sources with no identified counterpart at other wavelengths, and these unassociated sources have long been recognized as the natural testing ground for DM subhalo searches.

\subsection{Unassociated Sources in Fermi-LAT Catalogs}
\label{sec:5.4.1}

The latest incremental release, the 4FGL-DR4 catalog~\cite{2022ApJS..260...53A, 2023arXiv230712546B}, contains approximately 7,200 sources detected above a significance threshold of $4\sigma$.
Of these, 2,428 sources---roughly one third---remain unassociated, meaning that no plausible counterpart has been identified at other wavelengths~\cite{2023arXiv230712546B}.
Among the unassociated population, 1,282 sources lie at Galactic latitudes $|b| > 10{{}^\circ}$, where contamination from diffuse Galactic emission is reduced and the search for exotic source populations becomes more tractable.

Association completeness depends on both source brightness and Galactic latitude: brighter sources are far easier to link to multi-wavelength counterparts, while the association fraction drops from approximately 85\% at high latitudes to around 40\% near the Galactic plane, where dense source fields and X-ray absorption impede counterpart identification~\cite{2020ApJS..247...33A}.
The majority of unassociated sources are therefore expected to be ordinary astrophysical objects whose association has failed for instrumental or geometric reasons.
Their spectral properties offer direct clues to their nature: at high latitudes, the average photon index ($\Gamma \approx 2.24$) matches that of blazars of unknown type, while closer to the Galactic plane the spectra soften ($\Gamma \approx 2.42$), consistent with pulsars~\cite{2020ApJS..247...33A}.
Further discussion of the Fermi-LAT association procedure and its completeness can be found in Section~\ref{sec:2.3}.

To date, no unassociated Fermi-LAT source has been confirmed as a dark matter subhalo~\cite{2020ApJS..247...33A, 2022ApJS..260...53A}.
As discussed in Section~\ref{sec:5.3}, even under optimistic assumptions the most detectable subhalos would produce gamma-ray spectra similar to those of pulsars or soft-spectrum blazars, making individual identification on spectral grounds alone extremely difficult.
Nevertheless, the large unassociated population has motivated a sustained effort to search for DM subhalo candidates within it.

\subsection{Previous DM Subhalo Searches and the Failure of Classify-and-Count}
\label{sec:5.4.2}

All prior searches for dark matter subhalos among Fermi-LAT unassociated sources share a common strategy, referred to as \textit{classify-and-count}: a set of DM subhalo candidates is selected from the unassociated population based on their gamma-ray properties, and upper limits on $\langle\sigma v\rangle$ are derived under the assumption that the number of true subhalos cannot exceed the candidate count~\DMsubhaloUnas.
Early analyses relied on hand-crafted spectral and spatial cuts applied to successive catalog releases~\cite{2010PhRvD..82f3501B, 2012A&A...538A..93Z, 2012ApJ...747..121A, 2014PhRvD..89a6014B, 2015JCAP...12..035B, 2016JCAP...05..028S, 2017JCAP...04..018H}, with some studies emphasizing spatial extension as a potential discriminant~\cite{2014PhRvD..89a6014B, 2015JCAP...12..035B, 2022PhRvD.105h3006C}.
More recently, machine learning classifiers have automated and extended the candidate selection~\DMsubhaloML.
In all cases, none of the identified candidates could be unambiguously distinguished from standard astrophysical sources.

The classify-and-count approach, whether implemented with hand-crafted cuts or ML classifiers, suffers from three interrelated problems that compromise the statistical rigor of the resulting constraints.
First, no dark matter subhalo has ever been identified, so the ``DM'' class used in supervised classification has no empirical anchor and the resulting classification probabilities cannot be calibrated against reality~\cite{Amerio:2025fhz}.
Second, classifiers trained on balanced datasets with comparable fractions of AGNs, pulsars, and simulated DM sources bear no relation to the actual class composition of the gamma-ray sky, where the DM prevalence may well be zero; this constitutes a prior shift (Section~\ref{sec:3.4}) that systematically inflates the predicted DM probabilities~\cite{Amerio:2025fhz}.
Third, the number of DM candidates depends on an ad hoc probability threshold $p_\mathrm{DM} > p_\mathrm{threshold}$ whose value carries no statistical justification: different thresholds yield candidate counts that can differ by an order of magnitude and, consequently, different physical bounds on $\langle\sigma v\rangle$~\cite{Amerio:2025fhz}.
The resulting upper limits lack a well-defined confidence level---they represent an assumption ($N_\mathrm{DM} \leq N_\mathrm{candidates}$) rather than a statistical inference.

\subsection{The Quantification Learning Alternative}
\label{sec:5.4.3}

The failure of classify-and-count is conceptual, not technical: the approach attempts to answer a population-level question---how many DM subhalos contribute to the unassociated catalog?---using a tool designed for individual-level decisions.
A different framework is needed, one that models the composition of the unassociated population as a whole and fits the DM contribution directly to the data.

The formal machinery for this problem was introduced in Section~\ref{sec:3.4}: \textit{quantification learning} replaces the discriminative question ``what class does this source belong to?'' with the generative question ``what is the probability of observing features $\mathbf{x}$ given that the source belongs to class $k$?''
By estimating the class-conditional densities $p(\mathbf{x}|k)$ from the associated catalog and modeling the unassociated population as a mixture $p_\mathrm{target}(\mathbf{x}) = \sum_k \pi_k\, p(\mathbf{x}|k)$, the class prevalences $\pi_k$ become free parameters determined by the data rather than imposed during training~\cite{10.1145/3117807_Gonzalez_quantification, 2024arXiv240100490M}.
This generative formulation yields a well-defined product likelihood over all unassociated sources, provides a natural mechanism for including a new class (dark matter subhalos) whose prevalence was zero in the training set, and supports statistically rigorous confidence intervals through the profile likelihood~\cite{Amerio:2025fhz}.

The entanglement of prior and covariate shift between Fermi-LAT associated and unassociated populations---documented in Section~\ref{sec:3.4} and arising from the brightness-dependent association completeness---is handled simultaneously within this framework through sigmoid modulation functions that adapt the astrophysical templates to the unassociated feature space.
The full construction of this mixture model, including the kernel density estimation of astrophysical templates, the Monte Carlo simulation of the DM component, and the optimization via the Expectation-Maximization algorithm, is presented in the paper body that follows.
